\documentclass[9pt]{ltjsreport}
\usepackage{bm}
\usepackage{amsmath,amssymb}
\usepackage{tcolorbox}
\usepackage{graphicx}
\usepackage{booktabs}
\usepackage{subcaption}
\usepackage{tikz}
\usepackage{listings}
\usetikzlibrary{graphs}

\begin{document}

\section{線形計画}
\subsection{線形計画問題の標準形はどのような形か}
\begin{tcolorbox}[title=線形計画問題の標準形]
    \begin{gather*}
        \text{目的関数}:\bm{c}^\top \bm{x}
        \longrightarrow \min\\        
        \text{制約条件}:\bm{Ax}=\bm{b},\bm{x} \ge \bm{0}
    \end{gather*}
\end{tcolorbox}
\subsection{以下に示す問題系を標準形に変形せよ}

\begin{align*}
    \text{目的関数：}&-2x_1+5x_2\longrightarrow\max\\
    \text{制約条件：}&4x_1-6x_2=30\\
    &2x_1+8x_2\le50\\
    &7x_1+5x_2\ge10\\
    &x1\ge0,x_2\text{(符号成約なし)}
\end{align*}

\hrulefill

この問題は、標準形と比較し、
\begin{enumerate}
    \item 目的関数を最大化している
    \item $x_2$に符号制約がない
    \item $\le,\ge$を含む不等式がある
\end{enumerate}
という３点が異なる。

目的関数の最大化は、
目的関数全体に-1をかけた関数を最小化することと同じ。

(2)については、
変数$x_2$を
\begin{equation*}
    x_2=x_2'-x_2'', x_2'\ge0,x_2''\ge0
\end{equation*}
と置き換えることで対処できる。

以下では、$x_2'$を$x_2$、$x_2''$を$x_3$とかく。

(3)については、
$\le , \ge$のかわりに新しい非負変数(スラック変数)
を導入することで解決する。

以上をまとめると、次の問題は、
\begin{align*}
    \text{目的関数：}&2x_1-5x_2+5x_3\longrightarrow\mathrm{min}\\
    \text{制約条件：}&4x_1-6x_2+6x_3=30\\
                   &2x_1+8x_2-8x_3+x_4=50\\
                   &7x_1+5x_2-5x_3-x_5=10\\
                   &x_1,\dots,x_5\ge0
\end{align*}
と置き換えれる。

\newpage
\subsection{シンプレックス法とは}
前提として、線形計画問題の実行可能領域(制約条件を満たす\bm{x}の集合)
は空間$\mathbb{R}^n$内の凸多面体として存在する。

そして、最適解は一般に凸多面体の頂点の中にある。

２変数の線形計画問題を図示すればこれは容易にわかる。

全実行可能領域から最適解を見つけようとするのではなく、
その頂点を効率的に探索して最適解を見つけ出す手法が
シンプレックス法である。

\begin{tcolorbox}[title=基底解、実行可能基底解、基底変数、非基底変数、退化した基底解]
    線形計画問題において、その変数は$\bm{x}=(a_1,\dots,a_n)^\top$
    の形であると仮定する。そして、制約条件がk個あると仮定する。

    この時、自由度は\bm{x}の自由度はn-kである。

    n個の変数のうちn-k個の変数を適当に選び、それらに0を代入すれば
    残りの変数の値は一意的に定まるはずである。

    このようにして決まった$\bm{x}$を基底解という。

    基底解は$\binom{n}{n-k}$個存在する。
    \begin{description}
        \item[実行可能基底解] 基底解のうち$\bm{x}\ge\bm{0}$を満たすもの
        \item[非基底変数] 基底解を求める際に0とおいた変数
        \item[基底変数] 非基底変数以外の変数
        \item[退化した基底解] 基底変数の中に値が0であるようなものが存在する場合
    \end{description}
\end{tcolorbox}
\begin{tcolorbox}[title=基底行列、非基底行列]
    ある基底解\bm{x}は、n-k個の非底変数とk個の基底変数をを持つものとする。
    
    n-k次元非基底変数ベクトルを$\bm{x}_N$、
    k次元基底変数ベクトルを$\bm{x}_B$と置く。
    
    次に、
    制約条件$A\bm{x}=\bm{b}$を、基底変数に対応するk個の列と非基底変数
    に対応するn-k個の列に分割し、
    前者からなるAをB、後者をN、
    とおく。

    \begin{description}
        \item[基底行列]Bが正則である時のB
        \item[非基底行列]Bが正則であるときのN
    \end{description}
\end{tcolorbox}
制約条件式$A\bm{x}=\bm{b}$は基底行列と非基底行列で
\begin{gather}
    B\bm{x}_B+N\bm{x}_N=\bm{b}\\
    \bm{x}_B=B^{-1}\bm{b}, \bm{x}_N=\bm{0}
\end{gather}
と書ける。

$B^{-1}\bm{b}\ge0$の時、この基底解は実行可能基底解となる。

実行可能基底解はその各々が凸多面体の頂点に対応する。


\begin{tcolorbox}[title=ピポット操作]
ある実行可能基底解に対して、一組の基底変数と
非基底変数を入れ替えた基底解がまた実行可能基底解となる時、
この操作は実行可能領域の一つの頂点からそれと隣り合う
別の頂点に移動することを意味する。
このような基底変数と非基底変数の
入れ替えを\textbf{ピポット操作}という。
\end{tcolorbox}
\begin{tcolorbox}[title=シンプレックス乗数、双対コスト係数]
    以下のように定義されたm次元ベクトル$\bm{\pi}$をシンプレックス定数と呼ぶ。
    \begin{equation}
        \bm{\pi}=(B^\top)^{-1}\bm{c}_B
    \end{equation}

    以下のように定義されたベクトル
    \begin{equation}
    \bm{c}_N-N^\top\bm{\pi}
    \end{equation}
    を非規定変数$\bm{x}_N$の相対コスト係数と呼ぶ。
\end{tcolorbox}

\subsubsection{ある実行可能基底解が最適解かどうかの確認方法}
ある実行可能基底解$\bm{x}=(x_{B},x_{N})=(B^{-1}\bm{b},\bm{0})$を問題に代入する。
\begin{equation}
    \bm{x}_B=B^{-1}\bm{b}-B^{-1}N\bm{x}_N \label{eq:transformedBase}
\end{equation}
と表される。この式を代入し、問題を変形していくと、
線形計画問題は
\begin{gather}
    \text{目的関数：}\bm{\pi}^\top\bm{b}+
    (\bm{c}_N-N^\top\bm{\pi})^\top\bm{x}_N
    \longrightarrow \min\\
    \text{制約条件:}
    B^{-1}\bm{b}-B^{-1}N\bm{x}_N\ge\bm{0},
    \hspace{1cm}\bm{x}_N\ge\bm{0}    
\end{gather}
という形になる。

いま対象にしている以外の全ての実行可能解は$\bm{x}_N\ge\bm{0}$
であることを考えれば、
相対コスト係数が0以上
\begin{equation}
    \bm{c}_N-N^\top\bm{\pi}\ge 0
\end{equation}
であるならば対象の実行可能界は最適解であると言える。

\subsubsection{ある実行可能基底解が最適解でない時はどうするか}
    $\bm{c}_N-N^\top\bm{\pi}>0$を満たさない時、
    双対コスト係数の要素の中に負であるものが少なくとも
    一つ存在する。

    そのような負の係数をもつ非基底変数$x_k$を一つ選ぶ。
    この$x_k$を0から増加させれば、目的関数は減少する。

    もちろん、その増加は制約条件を満たしながら行わなければ
    ならないので、$\bm{x}_B$は式(\ref{eq:transformedBase})
    より、
    \begin{equation}
        \bm{x}_B=B^{-1}\bm{b}-B^{-1}\bm{a}_kx_k \label{eq:xBTrans}
    \end{equation}
    に従い変化する。($\bm{a}_k$はAのk列目を抜き出した列ベクトルとなる)
    (もとから$B^{-1}\bm{a}_kx_k$が差分として引かれた形)

    ここで、
    \begin{gather}
        \bar{\bm{b}}:=B^{-1}\bm{b}\\
        \bm{y}:=B^{-1}\bm{a}_k\\
    \end{gather}
    を定義し、式(\ref{eq:xBTrans})を書き換えれば、
    \begin{equation}
        x_i=\bar{b_i}-y_ix_k (i\in 基底変数の係数の集合)
    \end{equation}
    とかける。
    よって、非基底変数$x_k$を
    \begin{equation}
       \theta :=\min\{\frac{\bar{b_i}}{y_i}|y_i>0(i=1,2,\dots,m)\}
    \end{equation}
    まで増加させても、非負条件$\bm{x}\ge\bm{0}$は保たれる。

    $\theta$だけ$x_k$を増加させたなら、$\bm{x}_B$
    のなかに0となる値が現れる。

    このようにしてピポット操作を行う。

\begin{tcolorbox}[title=有界でない問題]
        仮に、$\forall (y_i<0)$であるなら、
    非基底変数$x_k$を任意に増加させたとしても、
    基底変数の非負条件を破ることはないため、
    目的関数を無限に小さくできる。
    このような問題を\textbf{有界でない}という。
\end{tcolorbox}

\subsubsection{シンプレックス法のまとめ}
以上に示したように、非基底変数と基底変数を入れ替えながら
実行可能解を探索することで最適解を見つける。

そのプロセスを以下にまとめる。
\begin{tcolorbox}[title=シンプレックス法のアルゴリズム(まとめ)]
\begin{enumerate}
    \item 線形計画問題を標準形にする
    \item 初期実行可能基底解として適当な基底解を一つ求める
    \item シンプレックス乗数を計算する
    \item 相対コスト係数を計算する
    \item 相対コスト係数が全て0以上なら、計算終了。
    そうでなければ、相対コスト変数のうち、負の係数をもつ
    非基底変数$x_k$を一つ選ぶ。
    \item $\bm{y}=B^{-1}\bm{a}_k$を計算する
    \item $\forall (y_i<0)$なら、問題は有界でないため終了。
    そうでないなら、\theta を計算し、\bm{x}を更新する。

    ステップ(3)に戻る。
\end{enumerate}
\end{tcolorbox}
実際に問題を解くときは、
以下で紹介するシンプレックスタブローを使って計算する。

\newpage

\subsection{シンプレックス・タブローとは}
シンプレックス・タブローとは
シンプレックス法の計算を表形式で簡単に行う方法である。

以下の線形計画問題をシンプレックスタブローで計算しろ。

\begin{align}
    \text{目的関数：}&-x_1-x_2\longrightarrow \min \\
    \text{制約条件：}&3x_1+2x_2+x_3=12\\
                   &x_1+2x_2+x_4=8\\
                   &x_1\ge 0,x_2\ge 0,x_3\ge 0,x_4\ge 0
\end{align}
この問題の初期実行可能基底解を$(0,0,12,8)$とする。

\hrulefill

上の問題に対し、以下のようなシンプレックス・タブローを考える。
\begin{table}[h]
        \centering
        \begin{tabular}{ccccc|c}\toprule
            &-1&-1&0&0&0\\ \midrule
            $x_3$&3&2&1&0&12\\
            $x_4$&1&2&0&1&8 \\ \bottomrule
        \end{tabular}
\end{table}

タブローの最初の行は目的関数における各変数の係数、
その下は制約条件の変数係数、右端の数は制約条件の右辺の定数を示す。

右上角にはこの基底解における目的関数値に-1をかけた数を書く。

左端の列には対応する基底変数を書く。

そして、このようにシンプレックスタブローを作れる場合、
双対コスト係数がタブローの最初の行に、
$\bm{x}_B$が一番右の列に現れることに注意する。

このタブローに対し、以下のような操作を行う。

\begin{enumerate}
    \item 双対コスト係数が全て正なら終了
    \item 双対コスト係数(一番上の行)が負となっている列のうち、
    その絶対値が最も大きい行を選択(第i列とする)
    \item 第i列の正の要素に注目し、その各要素でその行の
    右端の列の要素を割った値を計算(第i列が全て負の要素なら、
    この問題は有界でないので、解を持たない。)
    \item 計算した時、その値が最も小さくなる要素をピポット要素とする。
    \item ピポット要素で、列基本変形をする
    \item 基底変数（左端の列）を更新する
    \item 新しく得られたタブローで、１に戻る
\end{enumerate}
この操作は注意深く観察すれば、
先に述べたシンプレックス法のアルゴリズムと完全に
一致するとわかる。

この手順に沿って計算していくと、最適解が
$(x_1,x_2)^\top =(2,3)\top$
となる。

\newpage

\subsection{二段階シンプレックス法とは}
ある線形計画問題に対し、シンプレックス法を適用させようとしても、
問題に対し初期実行可能基底解が見つからないことがある。

このとき、初期実行可能解を見つけるために用いられるのが、
二段階シンプレックス法である。

以下の問題に対し二段階シンプレックス法を適用させよ。

\begin{align}
    \text{目的関数：}&-2x_1-x_2-x_3\longrightarrow \min \\
    \text{制約条件：}&x_1+2x_2=12\\
                   &x_1+4x_2+2x_3=20\\
                   &x_1\ge 0,x_2\ge 0,x_3\ge 0
\end{align}

\hrulefill

次のような補助問題を考える。

\begin{align}
    \text{目的関数：}&x_4+x_5\longrightarrow \min \\
    \text{制約条件：}&x_1+2x_2+x_4=12\\
                   &x_1+4x_2+2x_3+x_5=20\\
                   &x_1\ge 0,x_2\ge 0,x_3\ge 0,x_4\ge 0,x_5\ge 0
\end{align}
この時導入した$x_4,x_5$は人為変数、人工変数と呼ばれる。

仮に、補助問題において目的関数値が0となるような実行可能基底解
が得られた場合、$x_4=x_5=0$となるため、
$\bm{x}_b=(x_1,x_2,x_3)$は元の問題の制約条件も満たす
実行可能基底解となる。

補助問題における目的関数値が０となる実行可能基底解が存在すると
すれば、それは$x_1\ge 0,x_2\ge 0$から、そのまま最適解
となる。

このことから、補助問題の最適解を求め、その目的関数値が０だとすれば、
その最適解は元の問題の実行可能基底解となる。

まず、補助問題をタブローで表現しよう。

\begin{table}[h]
    \centering
    \begin{tabular}{cccccc|c}\toprule
        &0&0&0&1&1&0\\ \midrule
        $x_4$&1&2&0&1&0&12\\
        $x_5$&1&4&2&0&1&20 \\ \bottomrule
    \end{tabular}
\end{table}
上端の行は、双対コスト係数を表していない。

これを修正するには、
$x_4,x_5$に対応する行全体をそれぞれ一番上の行から引き、
\begin{table}[h]
    \centering
    \begin{tabular}{cccccc|c}\toprule
        &-2&-6&-2&0&0&-32\\ \midrule
        $x_4$&1&2&0&1&0&12\\
        $x_5$&1&4&2&0&1&20 \\ \bottomrule
    \end{tabular}
\end{table}
とすればいい。

このタブローに対し、計算すれば、
$(x_1,x_2,x_3,x_4,x_5)^\top =(4,4,0,0,0)^\top$
となる。

よって、主問題の実行可能基底解を$(x_1,x_2,x_3)=(4,4,0)$
とすればよい。

\newpage
\subsection{弱双対性定理、双対定理とは}
任意の線形計画問題に対し、もう一つの双対問題(Dual Problem)という
線形計画問題を定義できる。
この時、元の問題を主問題(Primal Problem)と呼ぶ。

標準形の問題に対する(D)は以下のように定義される。
\begin{tcolorbox}[title=標準形とその双対問題]
    \begin{gather}
        \text{目的関数}:\bm{c}^\top \bm{x}\label{eq:stddual}
        \longrightarrow \min\\        
        \text{制約条件}:A\bm{x}=\bm{b},\bm{x} \ge \bm{0}\notag
    \end{gather}
    \hrulefill
    \begin{gather*}
        \text{目的関数}:\bm{b}^\top \bm{w}
        \longrightarrow \max\\        
        \text{制約条件}:A^\top\bm{w}\le \bm{c}
    \end{gather*}
\end{tcolorbox}

対称性を考えた形では、以下のようになる。
\begin{tcolorbox}
    \begin{gather*}
        \text{目的関数}:\bm{c}^\top \bm{x}
        \longrightarrow \min\\        
        \text{制約条件}:A\bm{x}\ge \bm{b},\bm{x} \ge \bm{0}
    \end{gather*}
    \hrulefill
    \begin{gather*}
        \text{目的関数}:\bm{b}^\top \bm{w}
        \longrightarrow \max\\        
        \text{制約条件}:A^\top\bm{w}\le \bm{c},\bm{w}\ge 0
    \end{gather*}
\end{tcolorbox}

適当な一対の(P),(D)から、ほかの任意の(P),(D)の対を
導くことも可能である。

任意の(P),(D)の組に対して成立する以下の定理群を紹介する。
\begin{tcolorbox}[title=弱双対定理]
    (P)と(D)のそれぞれの任意の実行可能解$\bm{x},\bm{w}$
    に対して、常に不等式
    $\bm{c}^\top\bm{x}\ge\bm{b}^\top\bm{w}$
    が成立する。

    \hrulefill

    Proof.
    
    問題(\ref{eq:stddual})の主問題と双対問題
    の組で考える。

    それぞれの問題の制約条件から、

    \begin{equation*}
        \bm{c}^\top \bm{x}\ge (A^\top \bm{w})^\top \bm{x}
        =\bm{w}^\top\bm{b}
    \end{equation*}
    となる。
    
    \hrulefill

    Cor1.

    (P)の任意の実行可能解$\bm{x}$に対して、
    \begin{equation*}
        \bm{c}^\top\bm{x}\ge (D)\text{の最大値}
    \end{equation*}
    (D)の任意の実行可能解$\bm{w}$に対して、
    \begin{equation*}
        \bm{b}^\top\bm{w}\le(P)\text{の最小値}
    \end{equation*}    
    \hrulefill

    Cor2.

    (P)と(D)の実行可能解$\bm{x}$と$\bm{w}$が
    $\bm{c}^\top\bm{x}=\bm{b}^\top\bm{w}$
    を満たせば、$\bm{x}$と$\bm{w}$はそれぞれ
    (P)と(D)の最適解であり、
    (P)の最小値と(D)の最大値は等しい。

    \hrulefill

    Proof.

    弱双対性定理より、(P)の目的関数の値は常に
    (D)の目的関数の値より大きい。

    よって、
    $\bm{c}^\top\bm{x}=\bm{b}^\top\bm{w}$
    が成立するということは、(P)はこれ以上
    小さい解が存在せず、
    (D)はこれ以上大きい解が存在しないということである。

    このことから、Cor2が成り立つ。
    
    \hrulefill

    Cor3.

    (P)が有界でないならば、(D)は実行可能解をもたない。
    また、(D)が有界でないならば(P)は実行可能解を持たない。
\end{tcolorbox}

%clear
\begin{tcolorbox}[title=双対定理]
(P)または(D)の一方が最適解をもてば他方も最適解を持ち、
(P)の最小値と(D)の最大値は等しい

\hrulefill

Proof.

証明の方針として、(P)の最適解におけるシンプレックス乗数
が(D)においてどのような振る舞いをするかを探索する。

標準形である
問題(\ref{eq:stddual})の主問題とその双対問題の組で考える。

主問題が最適解を持つ時、
その双対問題も最適解を持つことを示せば十分である。

(P)の最適基底解を、
$\bm{x}^*_B=B^{-1}\bm{b},\bm{x}^*_N=\bm{0}$
とし、基底行列をB、非基底行列をNとかく。

ベクトル$\bm{w}^*$を$\bm{w}^*=(B^\top)^{-1}\bm{c}_B$
と定義する(シンプレックス乗数に同じ)。

するとBは最適基底行列であることから、双対係数より、
$\bm{c}_N-N^\top\bm{w}^*\ge\bm{0}$が成立する。

また、$\bm{w}^*$の定義から、$\bm{c}_B-B^\top\bm{w}^*=\bm{0}$
が成立する。これらから、

\begin{equation*}
\binom{B^\top}{N^\top}\bm{w}^*
\le\binom{\bm{c}_B}{\bm{c}_N}
\end{equation*}
の不等式が成り立つ。すなわち、$A^\top\bm{w}^*\le\bm{c}$
((D)の制約条件)を満たすので、$\bm{w}^*$は(D)の実行可能解
である。さらに、
\begin{equation*}
    \bm{c}^\top\bm{x}^*
    =\bm{c}_B^\top\bm{x}_B^*+\bm{c}_B^\top\bm{x}_N^*
    =\bm{c}_B^\top B^{-1}\bm{b}=(\bm{w}^*)^\top\bm{b}
\end{equation*}
であるから、弱双対性定理のCor2より、$\bm{x}^*$は(D)の最適解である。

\hrulefill

Cor.

(P)と(D)がともに実行可能解を持てば、それらは最適解をもち、
(P)の最小値と(D)
の最大値は等しい。

\hrulefill

Proof.

仮定から、(P)は実行可能解をもつ。

(D)は実行可能解を持つので、弱双対性定理のCor2より、
(P)は下に有界であることがわかる。

よって、(P)の最小値は存在する。すなわち、(P)は最適解
を持つので、双対定理より
(P),(D)は最適解をもち、かつ(P)の最小値と(D)の最大値
は等しくなる。

\end{tcolorbox}
これらの結果から、主問題と双対問題は実質的に等価な問題である
と示唆される。

\subsection{双対問題から主問題の最適解の求め方}
線形計画問題の解き方の一つに、その問題に対する双対問題
を利用する方法がある。

双対問題の最適解におけるシンプレックス乗数
をみることで、主問題の最適基底解
が見て取れる。

\newpage
\subsection{感度分析とは}
与えられた線形計画問題に対し、シンプレックス法で
最適解を計算できたとする。
現実においては、問題の係数に対し微妙な誤差が生じる
可能性が存在する。
では、その問題の係数の変化に対して最適解はどのように
振る舞うのだろうか。
このような考察を感度分析という。

標準形の問題において、制約条件の右辺の定数$\bm{b}=(b_1,
\dots,b_n)^\top$が$\bm{b}=(\Delta b_1,\dots,\Delta b_n)^\top$
だけ変化した問題を考える。
\begin{align}
    \text{目的関数}:\bm{c}^\top\bm{x}\longrightarrow \min\\
    \text{制約条件}:A\bm{x}=\bm{b}+\Delta\bm{b},\bm{x}\ge 0
\end{align}
$\Delta \bm{b}=\bm{0}$のときの最適基底解が
$\bm{x}_B=B^{-1}\bm{b},\bm{x}_N^*=\bm{0}$で既に計算されているものとする。

最適解の条件から、
\begin{gather}
    B^{-1}\bm{b}\ge\bm{0}\label{eq:kando}\\
    \bm{c}_N-N^\top(B^\top)^{-1}\bm{c}_B\ge\bm{0}
\end{gather}
を満たす。

$\bm{b}$が変化したとしても$(\bm{x}_B,\bm{x}_N)$
をが最適解であるためには、
\begin{equation}
    B^{-1}(\bm{b}+\Delta \bm{b})\ge\bm{0}
\end{equation}
を満たす必要がある。

このことから、仮に$\bm{x}_B$が非退化、すなわち
$\bm{x}_B^*=B^{-1}\bm{b}\ge\bm{0}$かつ
変化量が十分に小さいならば、最適基底行列Bは変化しない。

式(\ref{eq:kando})の目的関数の最小値を
$\varphi(\bm{b}+\Delta\bm{b})$とあらわせば、
\begin{gather}
    \varphi(\bm{b})
    =\bm{c}_B^\top B^{-1}(\bm{b})\\
    \varphi(\bm{b}+\Delta\bm{b})
    =\bm{c}_B^\top B^{-1}(\bm{b}+\Delta{b})\\
    \varphi(\bm{b}+\Delta\bm{b})-\varphi(\bm{b})
    =\bm{c}_B^\top B^{-1}(\Delta\bm{b})
\end{gather}
が成立する。

シンプレックス乗数が$\bm{w}^*=(B^\top)^{-1}\bm{c}_B$
となることを考えれば、
\begin{gather}
    \varphi(\bm{b}+\Delta\bm{b})-\varphi(\bm{b})
    =\bm{w}^{*\top}\Delta\bm{b}\\
    \frac{\partial\varphi(\bm{b})}{\partial b_i}=w_i^*　（i=1,\dots,m)\label{eq:partialShadowPrice}
\end{gather}
と書き換えることができる。

式(\ref{eq:partialShadowPrice})から、
制約条件の右辺を一単位変化させたときの目的変数の変化量は、
双対問題の解であることを示している。
\newpage
\subsection{シンプレックス法の計算量について述べよ}
シンプレックス法の反復回数は経験的に、問題の制約条件の
数の2から3倍である。
しかし、厳密に言うとシンプレックス法は多項式時間アルゴリズム
ではない。
ある種の人工的に作られた線形計画問題に対して、計算量
が爆発的に増加することがあり得る。
この欠点を克服したのがカーマーカーにより提唱された
内分法である。

内点法そのものは、領域の内部を経由して問題の解に収束する
点列を生成するという性質を持つアルゴリズムの総称である。
\subsection{パス追跡法を主双対内点法の枠組みで構成せよ}
\subsubsection{相補性条件による問題の書き換え}
標準形の形の(P)とその双対問題(D)を以下に再掲する。
\begin{tcolorbox}[title=標準形とその双対問題]
    \begin{gather}
        \text{目的関数}:\bm{c}^\top \bm{x}\label{eq:stddual}
        \longrightarrow \min\\        
        \text{制約条件}:A\bm{x}=\bm{b},\bm{x} \ge \bm{0}\notag
    \end{gather}
    \hrulefill
    \begin{gather*}
        \text{目的関数}:\bm{b}^\top \bm{w}
        \longrightarrow \max\\        
        \text{制約条件}:A^\top\bm{w}+\bm{s}= \bm{c},
        \bm{s}\ge \bm{0},\bm{w}\ge \bm{0}
    \end{gather*}
双対問題はスラック変数で書き換えられていることに注意
\end{tcolorbox}
双対定理、弱双対定理Cor2.より、
$\bm{x}$と$\bm{w}$が最適解であることと、
\begin{equation}
    \label{eq:naibunnPre}
    \begin{cases}
        \bm{c}^\top\bm{x}=\bm{b}^\top\bm{w}\\
        A\bm{x}=\bm{b}\\
        A^\top\bm{w}+\bm{s}=\bm{c}\\
        \bm{x}\ge\bm{0},\bm{s}\ge\bm{0}
    \end{cases}
\end{equation}
は等価である。

\begin{tcolorbox}[title=相補性条件]
(\ref{eq:naibunnPre})の第一式である
$\bm{c}^\top\bm{x}=\bm{b}^\top\bm{w}$
は、
\begin{equation}
    \bm{c}^\top\bm{x}-\bm{b}^\top\bm{w}
    =(A^\top+\bm{s})^\top\bm{x}-\bm{b}^\top\bm{w}
    =\bm{s}^\top\bm{x}+(A\bm{x}-\bm{b})^\top\bm{w}
    =\bm{s}^\top\bm{x}
\end{equation}
なので、$\bm{s}^\top\bm{x}=\bm{0}$
となるため、

X,Sをそれぞれ$x_j,s_j$を対角要素とする
$n\times n$行列、
$\bm{e}$を全ての要素が1である
n次元ベクトルとすると、
\textgt{相補性条件}
\footnote{(\ref{eq:naibunnPre})の第一式
を相補性条件に置き換えた、(\bm{x},\bm{w},\bm{s})
を求めるような問題は、その解を求めることで
元の問題の(P)、(D)の最適解を同時に求めることが
できるので、主双対最適性条件と呼ばれる。
この問題は、線形相補性問題というクラスに属する。
以下で説明する主双対内点法の考え方は、線形相補性問題
全体に対しても説明できる。}
\begin{equation}
    XS\bm{e}=\bm{0}
\end{equation}
で書き換えることが可能
\end{tcolorbox}
\newpage
\subsubsection{主双対内点法とは}
前節から、主問題、双対問題の組を
相補性条件で置き換えると、
\begin{equation}
    \label{eq:souhosei}
    \begin{cases}
        XS\bm{e}=\bm{0}\\
        A\bm{x}=\bm{b}\\
        A^\top\bm{w}+\bm{s}=\bm{c}\\
        \bm{x}\ge\bm{0},\bm{s}\ge\bm{0}
    \end{cases}
\end{equation}
となる。
この問題に対し、
新たにパラメータ$\mu>0$を定義し
\begin{equation}
    \label{eq:paramMu}
    \begin{cases}
        XS\bm{e}=\mu\bm{e}\\
        A\bm{x}=\bm{b}\\
        A^\top\bm{w}+\bm{s}=\bm{c}\\
        \bm{x}>\bm{0},\bm{s}>\bm{0}
    \end{cases}
\end{equation}
を考える。

$\mu\longrightarrow 0$とすれば、
\ref{eq:souhosei}の解に収束する。

そこで、0に収束する数列$\{\mu^k\}_{k\in N}$
を考え、これが\ref{eq:paramMu}の$\mu$
パラメータに入るとする。


このとき、$\mu^{k-1}$の時の近似解が得られている
のであれば、
\ref{eq:paramMu}を線形近似した方程式で、
$\mu^k$の時の近似解を
$\mu^{k-1}$の時の近似解の差分として求めることが可能。

このようにして近似を求めていくと、それは
元の問題の解に収束することが知られている。
\end{document}